% ============================================================
\section{The Horizon Fundamental Theorem}
\label{sec:hft}
% ============================================================

This section records the structural projection laws behind the calculus.
The order matters:

\begin{enumerate}[label=(\roman*)]
\item for any square-zero additive operator, projection produces an exact defect identity;
\item for any compatible horizon tower, orthogonal refinement produces an exact telescoping law in energy;
\item realized models of the calculus may then instantiate these structural laws with canonical concrete data.
\end{enumerate}

The canonical closure-balance realization of these laws is deferred to
Part~I. This chapter develops only the structural theorem layer.

\begin{theorem}[General HFT-1]
\label{thm:HFT1-general}\lean{HFT1_general_nilpotent}
Let $\cH$ be an additive state space with a transitive horizon tower $(\CCPi{h})_{h\in\bbN}$, let $A$ be an additive operator on $\cH$, and assume
\[
A^2 = 0.
\]
Define the observed operator $\CCObs{A}{h} := \CCPi{h} A \CCPi{h}$ and $\CCPibar{h} := I - \CCPi{h}$.
Then
\[
\CCObs{A}{h}^2 \;=\; -\,\CCPi{h} A \CCPibar{h} A \CCPi{h}.
\]
Equivalently,
\[
\CCObs{A}{h}^2 \;=\; -\,\CCDefectExpr{A}{h}.
\]
\end{theorem}

\begin{proof}
Compute
\[
0 = \CCPi{h} A^2 \CCPi{h}
= \CCPi{h} A(\CCPi{h} + \CCPibar{h})A\CCPi{h}
= \underbrace{\CCPi{h} A\CCPi{h} A\CCPi{h}}_{\CCObs{A}{h}^2} + \CCPi{h} A \CCPibar{h} A \CCPi{h}.
\]
Rearranging gives the claim.
\end{proof}

The remaining part of the Horizon Fundamental Theorem is the refinement law.
Its algebraic content is that adjacent refinement layers close internally and telescope.
Its metric content is that, when those layers are orthogonal, the hidden remainder decomposes into an exact sum of layer energies.

\begin{theorem}[General HFT-2]
\label{thm:HFT2-general}\label{thm:HFT}\lean{HFT2_finite_energy}\lean{HFT2_faithful_stabilized}
Let $(\CCPi{h})_{h\in\bbN}$ be a nested horizon tower on $\cH$ whose refinement layers split orthogonally, so that each horizon cut satisfies an exact Pythagorean decomposition.
For nested horizons $h_0 \le h_1 \le \cdots \le h_N$ and any $x\in\cH$,
\[
\|x - \CCPi{h_0} x\|^2
=
\sum_{j=0}^{N-1} \|(\CCPi{h_{j+1}} - \CCPi{h_j})x\|^2 + \|x - \CCPi{h_N} x\|^2.
\]
If the tower is \emph{tail-faithful in the HFT-2 sense}, meaning that the tail
energy eventually vanishes along consecutive refinement, then there exists $N$
such that for every $k \ge 0$,
\[
\|x - \CCPi{h_0} x\|^2
=
\sum_{j=0}^{N+k-1} \|(\CCPi{h_0+j+1} - \CCPi{h_0+j})x\|^2.
\]
\end{theorem}

\begin{proof}
The finite statement is the orthogonal telescoping law for the refinement chain: adjacent internal interfaces cancel algebraically, and orthogonality converts that telescoping decomposition into an exact sum of energies.
The tail-faithful statement is the exact discrete stabilization of those finite
prefix sums once the tail energy vanishes.
No analytic infinite-series machinery is used at this stage.
\end{proof}

\begin{remark}[HFT-2 as realized exactification ledger]
\label{rem:hft2-realized-exactification-ledger}\lean{CoherenceCalculus.TerminationData, CoherenceCalculus.nonterminal_incidences_continue, CoherenceCalculus.pair_telescope_energy}
At the bedrock level, the termination law identifies zero defect with
terminality, while nonterminal incidence continues by strict extension and
closure defect can vanish only on strict extension if terminality is
eventually reached
(Axioms~\ref{ax:termination} and~\ref{ax:nonretroactive-exactification}).
HFT-2 is the first realized telescoping law for that continuation burden:
adjacent refinement layers cancel internally, and the remaining quantity is an
exact sum of realized layer energies. In this sense HFT-2 does not generate
continuation; it measures the additive burden of exactification once a horizon
tower has been chosen.
\end{remark}

\begin{remark}[Tail faithfulness versus continuum-bridge faithfulness]
\label{rem:hft2-tail-faithfulness}
\lean{HFT2_faithful_prefix_exact, HFT2_faithful_stabilized}
The HFT-2 use of ``tail-faithful'' is a finite-stabilization condition: after
some finite refinement stage, the hidden tail energy vanishes exactly.
This is stronger than, and distinct from, the continuum-bridge notion of
faithfulness as strong limit recovery along an infinite tower.
\end{remark}

\begin{lemma}[One-step shadow-energy split]
\label{lem:hft2-one-step}\lean{pair_telescope_energy}
Under the hypotheses of Theorem~\ref{thm:HFT2-general}, every nested pair of horizons $h_0 \le h_1$ satisfies
\[
\|x - \CCPi{h_0} x\|^2
=
\|(\CCPi{h_1} - \CCPi{h_0})x\|^2 + \|x - \CCPi{h_1} x\|^2.
\]
\end{lemma}

\begin{proof}
Split the $h_0$-shadow into the part revealed at horizon $h_1$ and the remaining $h_1$-shadow.
Orthogonality of the refinement split gives the exact Pythagorean identity.
\end{proof}

\begin{corollary}[Three-level HFT-2]
\label{cor:hft2-three-level}\lean{HFT2_three_level_energy}
Under the hypotheses of Theorem~\ref{thm:HFT2-general}, every triple $h_0 \le h_1 \le h_2$ satisfies
\[
\|x - \CCPi{h_0} x\|^2
=
\|(\CCPi{h_1} - \CCPi{h_0})x\|^2
+ \|(\CCPi{h_2} - \CCPi{h_1})x\|^2
+ \|x - \CCPi{h_2} x\|^2.
\]
\end{corollary}

\begin{proof}
Apply Lemma~\ref{lem:hft2-one-step} first to $(h_0,h_1)$ and then to $(h_1,h_2)$.
\end{proof}

\begin{proposition}[Faithful exact-prefix capture]
\label{prop:hft2-prefix-exact}\lean{HFT2_faithful_prefix_exact}
If the tower is tail-faithful in the sense of
Theorem~\ref{thm:HFT2-general}, then for every initial horizon $h_0$ there
exists a finite stage $N$ such that
\[
\|x - \CCPi{h_0} x\|^2
=
\sum_{j=0}^{N-1} \|(\CCPi{h_0+j+1} - \CCPi{h_0+j})x\|^2.
\]
So the total hidden energy is already captured by a finite prefix of consecutive refinement increments.
\end{proposition}

\begin{proof}
By tail faithfulness, the tail energy vanishes after sufficiently many
consecutive refinements.
At that stage the finite telescoping identity has zero remainder, so only the prefix sum remains.
\end{proof}

\begin{proposition}[Structured infinite choice along a refinement tower]
\label{prop:structured-infinite-choice}\lean{RootedFiniteRefinementTower}\lean{structured_infinite_choice}
Let $(X_n)_{n\in\bbN}$ be a rooted refinement tower with predecessor maps
\[
p_n : X_{n+1} \to X_n,
\]
such that:
\begin{enumerate}[label=(\roman*),itemsep=0.2em]
\item the initial stage $X_0$ is given by an explicit nonempty finite list of roots;
\item for each $x \in X_n$, the refinement fiber above $x$ is given by an explicit nonempty finite list;
\item every listed refinement projects back to its parent under $p_n$.
\end{enumerate}
Then there exists a coherent infinite thread $(x_n)_{n\in\bbN}$ with
\[
x_n \in X_n,
\qquad
p_n(x_{n+1}) = x_n
\quad\text{for all }n.
\]
The construction is explicit: choose the first root, then choose the first listed refinement above it at every stage.
\end{proposition}

\begin{proof}
Start from the first listed root in $X_0$. Given $x_n$, take the first listed refinement above it and call that $x_{n+1}$. By hypothesis each refinement fiber is nonempty, so the recursion never stops, and by hypothesis each chosen refinement projects back to its parent. This gives an infinite coherent thread.
\end{proof}

\begin{quote}\small
\textbf{Scope.} This is not the full Axiom of Choice. It is a constructive infinite-choice principle for explicitly finite refinement systems. Its role here is to show how the finite tower grammar already supports coherent infinite selection in the structured limit.
\end{quote}

\begin{corollary}[Countable structured choice from finite partial choices]
\label{cor:countable-structured-choice}\lean{ExplicitFiniteFamily}\lean{countable_structured_choice}
Let $(F_n)_{n\in\bbN}$ be a countable family of explicitly finite nonempty fibers. For each finite stage $n$, let a \emph{partial choice} mean a compatible choice on the first $n$ fibers. Then there exists:
\begin{enumerate}[label=(\roman*),itemsep=0.2em]
\item a global sequence $(c_n)_{n\in\bbN}$ with $c_n \in F_n$ for every $n$;
\item a coherent chain of partial choices whose $(n+1)$st stage is obtained from its $n$th stage by adjoining $c_n$.
\end{enumerate}
\end{corollary}

\begin{proof}
The stage-$n$ partial choices form a rooted finite refinement tower: the predecessor map drops the last coordinate, and every stage extends by choosing one element from the next explicit finite fiber. Proposition~\ref{prop:structured-infinite-choice} therefore yields a coherent infinite thread of partial choices. Reading the new coordinate at each stage gives the global sequence $(c_n)$.
\end{proof}

\subsection{Dyadic refinement bridge}

The constructive refinement principles above already support a first continuum-like layer without importing reals, quotients, or field structure.
The bridge is explicit and dyadic: binary choices build a tower of adjacent endpoint data on finer and finer unit meshes.

\begin{proposition}[Finite dyadic address injectivity]\label{prop:dyadic-address-injective}\lean{partialChoice_snoc_eq_iff}\lean{refineDyadicCell_eq_iff}
Extending a finite dyadic address by one more binary digit is injective in both pieces of explicit data.
Equivalently, if two one-step refinements agree, then their parent addresses agree and their final digits agree.
\end{proposition}

\begin{proof}
At the finite level, a dyadic address is built by adjoining one last binary digit to a shorter address.
So equality of two refined addresses forces equality of the shorter addresses and equality of the final digits separately.
\end{proof}

\begin{proposition}[Dyadic interval refinement]\label{prop:dyadic-refinement}\lean{dyadicCell_two_children}\lean{child_intervals_distinct}\lean{intervalsOfDigits_leftIndex_step}\lean{intervalsOfDigits_rightIndex_step}\lean{intervalsOfDigits_mesh}
At depth $n$, binary refinement determines dyadic interval data with adjacent endpoint indices $(\ell_n,r_n)$ satisfying
\[
r_n = \ell_n + 1,
\qquad
r_n \le 2^n.
\]
Each stage has exactly two distinct children.
If the next digit is left, then
\[
\ell_{n+1} = 2\ell_n,
\qquad
r_{n+1} = 2\ell_n + 1.
\]
If the next digit is right, then
\[
\ell_{n+1} = 2\ell_n + 1,
\qquad
r_{n+1} = 2r_n.
\]
\end{proposition}

\begin{proof}
This is the direct recursion of binary refinement.
Left refinement takes the left half of the current dyadic cell, while right refinement takes the right half; the endpoint formulas are exactly the resulting doubled-index updates on the finer mesh.
\end{proof}

\begin{proposition}[Binary threads yield dyadic point-candidates]\label{prop:dyadic-point-candidates}\lean{binarySequence_exists}\lean{unitPointOfDigits}\lean{leftUnitPoint}\lean{intervalsOfDigits_nested_scaled}\lean{unitPoint_mesh_strict}
Every infinite binary digit sequence determines a coherent dyadic interval thread on the unit mesh.
At each step, the next interval is nested inside the previous one after passing to the common finer denominator, and the mesh strictly refines from $2^n$ to $2^{n+1}$.
Hence the present structural layer already supports quotient-free dyadic
unit-point candidates.
\end{proposition}

\begin{proof}
Read each digit as a left/right child choice in the binary refinement tower.
The resulting endpoint recursion gives a coherent nested interval thread, and the denominator doubles at every step.
Taking the constant-left sequence gives a canonical example.
\end{proof}

\begin{quote}\small
\textbf{Scope.} This is still a pre-real layer.
Ambiguous binary tails have not yet been identified, so no quotient of equivalent dyadic expansions is taken here.
The present structural layer already forces a rigorous shrinking-thread
scaffold above the discrete level.
\end{quote}

\begin{quote}\small
\textbf{Intuition (HFT-1).} The full boundary closes: $d^2=0$ on the uncropped coherence space. But the \emph{observed} boundary $\CCObs{d}{h}$ is the cut $\CCPi{h} d \CCPi{h}$, and that cut need not close. The failure is exactly the ``round trip through the shadow'': apply $d$, some escapes to the shadow, apply $d$ again, some returns. That returned part is the defect.
\end{quote}

\begin{quote}\small
\textbf{Intuition (HFT-2).} The total ``distance'' from $x$ to its coarsest projection decomposes into orthogonal increments at each refinement level. This is the Pythagorean theorem applied to a tower of nested subspaces, with integration expressed as a sum of layer-by-layer contributions.
\end{quote}

\begin{tcolorbox}[colback=yellow!5, colframe=black!70, title={Warning: Zero defect $\neq$ conservation}]
\label{box:zero-defect-warning}\lean{returnOp}\lean{defect_factorization}
\textbf{The three operators.} For horizon projection $\CCPi{h}$ with $\CCPibar{h} = I - \CCPi{h}$:
\begin{align*}
L_h &:= \CCPibar{h} d \CCPi{h} &&\text{(leakage: observable $\to$ shadow)}, \\
R_h &:= \CCPi{h} d \CCPibar{h} &&\text{(return: shadow $\to$ observable)}, \\
\CCDefect{d}{h} &:= R_h L_h = \CCPi{h} d \CCPibar{h} d \CCPi{h} &&\text{(defect: round-trip)}.
\end{align*}

\textbf{Key distinctions:}
\begin{itemize}[itemsep=0.2em]
\item $\CCDefect{d}{h} = 0$ is \emph{algebraic closure} (a compositional identity), \textbf{not} metric conservation.
\item $L_h \neq 0$ with $\CCDefect{d}{h} = 0$ is possible when $R_h = 0$ (``silent leak'').
\item Orthogonality of $\CCPi{h}$ is what makes the defect energy $\|L_h x\|^2$ well-defined.
\end{itemize}
\end{tcolorbox}

\begin{lemma}[Silent leak configuration]
\label{lem:silent-leak}\lean{silent_leak_configuration}
There exist $\cH$, a projection $\Pi$, and $d$ with $d^2 = 0$ such that:
\begin{enumerate}[label=(\roman*),itemsep=0.1em]
\item $d_\Pi := \Pi d \Pi$ satisfies $d_\Pi^2 = 0$ (algebraic closure),
\item $L := (I - \Pi) d \Pi \neq 0$ (nonzero leakage),
\item $\delta := \Pi d (I - \Pi) d \Pi = 0$ (zero defect, because $R = \Pi d (I-\Pi) = 0$).
\end{enumerate}
\end{lemma}

(Proof in Appendix~\ref{sec:appendix-tower}.)

\subsection{Illustration: The 4-vertex path}

Returning to the running example (Example~\ref{ex:hello-world}): the
4-vertex path with block-averaging horizon. HFT-1 applies because the graded
differential $D$ satisfies $D_{k+1} \circ D_k = 0$
(Definition~\ref{def:nilpotent-differential}). The defect term is exactly the
within-block contribution that the coarse observer cannot represent. This
subsection remains illustrative; the theorem-bearing statements are the general
and coherence-realized HFT identities above.

\subsection{Defect operator and defect energy}

\begin{definition}[Defect operator]\label{def:defect-operator}\lean{defectOp}
Define the \emph{defect operator}
\[
\CCDefect{d}{h} := \CCDefectExpr{d}{h}
\quad\text{so that}\quad
\CCObs{d}{h}^2 = -\CCDefect{d}{h}.
\]
\end{definition}

\begin{definition}[Defect energy]\label{def:defect-energy-operator}\lean{defectEnergy}\lean{defectEnergy_operator_form}
Define the \emph{defect energy} on $\cH$ by
\[
\CCDefEnergy{h}{x} := \|\CCLeak{d}{h} x\|^2 = \|\CCLeakExpr{d}{h}\, x\|^2.
\]
Whenever the leakage map admits an energy adjoint $\CCLeak{d}{h}^\ast$ for the coordinate pairing of the state space, this can be written in operator form as
\[
\CCDefEnergy{h}{x} = \langle x, \CCDefOp{d}{h} x\rangle,
\qquad
\CCDefOp{d}{h} := \CCLeak{d}{h}^\ast \CCLeak{d}{h}.
\]
\end{definition}

\paragraph{Why this is a calculus.}
The defect $\CCDefect{d}{h}$ is the derivative object: it measures failure of
closure under observation. The telescoping identity (HFT-2) is the integration
analog: total deviation decomposes into orthogonal layer increments. Defect
energy $\CCDefEnergy{h}{x} = \|\CCLeak{d}{h} x\|^2$ quantifies the portion of
$dx$ that lands in the shadow and is therefore invisible at horizon $h$.

\paragraph{Interpretation note.}
\label{rem:interpretation-note}
The defect energy is an algebraic quantity measuring representational mismatch,
not physical loss. Finer horizons do not make defect energy monotone for
arbitrary states; that requires additional representability hypotheses, as made
explicit later in Proposition~\ref{prop:log-defect}. In a faithful tower, the
tail term in HFT-2 eventually vanishes after finite stabilization. No
thermodynamic or information-theoretic interpretation is claimed here without
additional structure.

The next section records chain and product rules for coherence defects under composition, analogous to the Leibniz and chain rules of classical calculus.
